% Beta-minus decay as Fermi's four-point contact interaction (the "blob"),
% the left half of fermi-vs-W.tex on its own. Uses feynmp via feynmp-auto.
% Build (in this directory):
%   pdflatex -shell-escape fermi-blob.tex && pdflatex -shell-escape fermi-blob.tex
%   pdftocairo -svg fermi-blob.pdf ../fermi-blob.svg
\documentclass[border=6pt]{standalone}
\usepackage{feynmp-auto}
\usepackage{xcolor}

\newif\ifdark \darktrue          % \darkfalse = black on white
\newcommand{\blobsize}{4mm}      % size of the four-point blob
\newcommand{\figfont}{\LARGE}    % font size for labels

\ifdark \pagecolor{black} \fi
\unitlength=1mm

\begin{document}
\ifdark \color{white} \fi
\figfont
\begin{fmffile}{fermiblob}
  \ifdark \fmfcmd{foreground := white; background := black;} \fi
  \fmfset{thin}{0.8pt}
  \fmfset{arrow_len}{3mm}
  \fmfframe(4,8)(4,8){% margins (left,top)(right,bottom) leave room for labels
  \begin{fmfgraph*}(36,30)
    \fmfleft{n,p}
    \fmfright{nub,e}
    \fmf{fermion}{n,v}
    \fmf{fermion}{v,p}
    \fmf{fermion}{v,e}
    \fmf{fermion}{nub,v}
    \fmfv{decor.shape=circle,decor.filled=full,decor.size=\blobsize}{v}
    \fmflabel{$n$}{n}
    \fmflabel{$p$}{p}
    \fmflabel{$e^-$}{e}
    \fmflabel{$\bar\nu_e$}{nub}
  \end{fmfgraph*}}
\end{fmffile}
\end{document}
